\section{Discussion}
\subsection{What ``structural versus heuristic'' means}
The central result is not that heuristics are useless. In-distribution, the fuzzy controller challenges 88\% of suspicious prompts, and such triage can support step-up authentication, logging, or human review. The result is that request-derived suspicion features are a poor foundation for a confidentiality boundary when modest reformulation removes the feature signal. In our adaptive set, the attacker forces lexical injection risk to zero in every case while preserving disclosure intent. A control whose security depends on that feature therefore loses most of its recognition capability.

The ACL behaves differently because its input is authenticated policy state. A paraphrase can change the words, but not whether the target belongs to the caller's authorized set. In the modeled path, \sbsr{}=1.0 is therefore guaranteed by candidate-set construction provided that identity, policy state, and implementation are correct; the benchmark confirms that implementation correspondence for the tested inputs. The design implication is precise: \emph{heuristic layers may narrow, challenge, or monitor; they should not be the mechanism that widens or defines the sensitive retrieval boundary}.

\subsection{Context exposure versus generated disclosure}
Output-only evaluation can dramatically understate exposure. Smol prompt-only \udr{} is 0.286\% while \ucer{} is 100\%; Qwen \udr{} is higher but still far below \ucer{}. A model refusal does not retroactively remove data already placed in its context, and future model/library changes can alter disclosure behavior. Security evaluation should therefore report both what the retriever exposed and what the generator emitted.

\subsection{Scale interpretation}
The scale experiment also illustrates why security/system claims require careful controls. ``100K RAG'' can mean a 100K global corpus with a fixed 200-record authorization domain, a 20K proportional domain, or an unfiltered 100K domain. Those are different workloads. Our results do not show that authorization becomes more secure at larger scale; they show that policy scope controls how much of global growth reaches retrieval. This distinction should be explicit in future evaluations.

\section{Cross-Model and Mechanistic Interpretation}
The two-model comparison helps separate claims that belong to the architecture from those that belong to generation behavior. Prompt-only \ucer{} is 100\% for both models because the global retriever presents unauthorized targets irrespective of how the model later responds. Secured \ucer{} is 0\% for both because candidate construction excludes out-of-scope records. These outcomes are therefore stable across model substitution. By contrast, prompt-only \udr{} differs by more than an order of magnitude between SmolLM2 (0.286\%) and Qwen2.5 (3.086\%) on the primary unauthorized set. This is exactly the kind of quantity that should not be elevated into an architectural guarantee: the context is the same type of unauthorized exposure, but the model's tendency to reproduce a canary differs substantially.

Utility is similarly model-dependent. At the original $k=2$, Qwen's prompt-only, pre-ACL, and risk-aware \arsr{} values are 92.5\%, 90.5\%, and 91.5\%, a much smaller spread than Smol's. With $k=1$, each model produces identical \arsr{} across variants: 89.5\% for SmolLM2 and 92.0\% for Qwen. Because exact target context, model, and greedy decoding are then identical, within-model equality is expected by construction; the result functions as a sanity check isolating retrieval composition rather than as a surprising utility discovery.

The heuristic result has a different structure. The fuzzy and deterministic controllers have different primary operating points, but both lose most of their suspicious-request challenge capability on shifted wording. Fuzzy \hdr{} of approximately 0.19--0.20 means that only about one fifth of the controller's original challenge propensity survives. The adaptive mutation search strengthens the interpretation because it directly optimizes the declared lexical feature rather than relying only on human-written paraphrases. Nevertheless, the attacker is intentionally constrained; a stronger semantic search could discover different weaknesses, and the observed retention ratio should not be generalized beyond this detector family without replication.

Ablation closes the explanatory loop. If injection-risk recognition were itself the confidentiality boundary, driving it to zero would be expected to open unauthorized retrieval. It does not: when authorization confidence remains intact, unauthorized deny stays at 100\% on shifted sets. Conversely, neutralizing authorization confidence eliminates shifted unauthorized denial despite leaving the other heuristic features present. The experiment therefore provides a mechanistic decomposition: request-derived features influence suspicion treatment, whereas policy-derived eligibility determines sensitive-record reachability in the implemented pipeline.

This decomposition is useful beyond a binary secure/insecure comparison. It suggests that RAG evaluations should classify controls by what state they trust and where they execute. A classifier or fuzzy score can be valuable for anomaly response, but its decision surface is exposed to language variation. An ACL can also fail, but its failure modes are different: stale entitlements, identity compromise, policy bugs, or metadata corruption rather than lexical paraphrase. Security engineering benefits when those mechanisms are tested against threats appropriate to their actual dependency surfaces instead of averaging them into a single guardrail score.

\section{Limitations and Threats to Validity}
The corpus is synthetic and the ACL is intentionally simple. Synthea avoids real-patient exposure but cannot reproduce all semantics, access relationships, or operational pressures of a clinical deployment. We test two small open instruction models, not frontier proprietary models or a broad model suite. The structural result is largely model-independent because it concerns which records enter context, but generation-level disclosure and utility remain model-dependent, as the Smol/Qwen \udr{} difference demonstrates.

The adaptive attacker is white-box: it optimizes a finite, source-controlled mutation neighborhood by directly calling the fixed 22-entry lexical detector's scoring functions, and its vocabulary was designed with knowledge of that detector. The result must not be generalized to learned, embedding-based, LLM-based, or black-box detectors. Semantic preservation uses automated encoder and LLM checks rather than blinded human annotation; moreover, Qwen is both the 19-template semantic judge and one subject model, creating a limited model-family overlap. A future study should add independent annotators and genuinely black-box semantic adversaries.

\sbsr{}=1.0 follows from the modeled candidate-filtering construction when its assumptions hold and is empirically confirmed on the tested path; it is not proof of zero real-world leakage. Side channels, identity compromise, authorization-data errors, poisoned metadata, vector-store bugs, or indirect context flows are outside scope. Similarly, cryptographic confidentiality and formally verified implementation are not provided; provably secure/encrypted RAG addresses stronger guarantees \cite{zhou2025provable}.

The scale experiment uses a specific TF-IDF retrieval representation and 90 fixed legitimate generation queries per cell. The fixed-scope flatness is intentionally a construction control. Differences of one or a few successes should not be interpreted as broad ranking of modes without larger datasets and multiple retrieval families.

\section{Implications for Secure RAG Design}
The evidence supports five practical principles. First, enforce record eligibility before sensitive retrieval/context assembly. Second, make heuristic risk scoring subordinate: it may reduce access, trigger confirmation, or increase auditing, but must not expand the policy-visible corpus. Third, evaluate context exposure separately from output leakage. Fourth, report security and utility at the cluster/template level when prompts are systematic variations of fewer attack strategies. Fifth, distinguish global corpus size, absolute authorization scope, and authorization density when discussing scalability.

These principles align with recent pipeline-oriented RAG security surveys \cite{xu2026taxonomy,palanisamy2026survey} while adding an empirical demonstration of why defense layers should be evaluated by mechanism rather than grouped into one guardrail score. The framework may generalize beyond healthcare to any RAG deployment where identity/policy state and behavior-derived risk signals coexist, although that external validity remains to be tested.

\section{Conclusion}
We evaluated a simple but consequential distinction in RAG security. Structural authorization and heuristic risk triage can coexist in the same pipeline, yet they generalize differently. Across two language models and primary, held-out, and adaptive adversarial distributions, pre-retrieval ACL enforcement maintained the tested authorization boundary, while heuristic suspicious-request recognition retained only about one fifth of its primary challenge rate after reformulation. Feature ablation shows why: policy-derived authorization confidence drives unauthorized denial under shift, whereas lexical injection-risk can be adversarially driven to zero. A second model reproduces the controlled utility result, and a fresh four-mode scale matrix shows that caller-visible authorization scope---not just global corpus size---determines retrieval workload. The broader lesson is not that heuristics should be removed, but that sensitive context reachability should not depend on their success. In secure RAG systems, heuristics can recognize risk; structural controls should constrain capability.
